\addcontentsline{toc}{subsubsection}{Solution 46. The Missing Middle Value}

\begin{dlsummary}[title={\bfseries Solution 46. The Missing Middle Value},breakable]

\begin{proof}
The only solutions are
\[
\boxed{f(n)=n\qquad\text{and}\qquad f(n)=n^2.}
\]

\paragraph{1. Establish order and bound \(f(2)\).}

Multiplicativity and positivity give \(f(1)=1\). Superadditivity then implies
\[
f(n+1)\ge f(n)+1.
\]
Thus \(f\) is strictly increasing and \(f(n)\ge n\).

Write \(k=f(2)\), so \(k\ge2\). If \(k\ge5\), strict increase and multiplicativity give
\[
f(f(2))=f(k)>f(4)=k^2\ge25,
\]
contradicting \(f(f(2))\le16\). Therefore
\[
\boxed{k\in\{2,3,4\}.}
\]

\paragraph{2. Eliminate the middle value.}

Suppose \(f(2)=3\). Since \(9>8\),
\[
f(3)^2=f(9)>f(8)=27.
\]
Because \(f(3)\) is an integer, this forces \(f(3)\ge6\).

But \(243<256\), so
\[
f(3)^5=f(243)<f(256)=3^8.
\]
This is impossible, since
\[
f(3)^5\ge6^5=7776>6561=3^8.
\]
Hence
\[
f(2)=2\quad\text{or}\quad f(2)=4.
\]

\paragraph{3. One value determines the whole function.}

Write \(f(2)=2^d\), where \(d\in\{1,2\}\). Fix \(n\ge2\).

For each positive integer \(r\), choose an integer \(s\ge0\) such that
\[
2^s\le n^r<2^{s+1}.
\]
Monotonicity and multiplicativity give
\[
2^{ds}\le f(n)^r<2^{d(s+1)}.
\]
Also,
\[
2^{ds}\le n^{dr}<2^{d(s+1)}.
\]
Dividing the bounds appropriately,
\[
2^{-d}<\left(\frac{f(n)}{n^d}\right)^r<2^d.
\]
These bounds hold for every positive integer \(r\). A positive number other than \(1\) has powers tending either to \(0\) or to infinity, so
\[
f(n)=n^d.
\]
Together with \(f(1)=1\), this proves the classification.

\paragraph{4. Verify both solutions.}

For \(f(n)=n\), all conditions are immediate. For \(f(n)=n^2\),
\[
(m+n)^2\ge m^2+n^2,
\qquad
f(f(n))=n^4,
\]
and multiplicativity also holds.
\end{proof}

\paragraph{Takeaways.}
\begin{itemize}
\item **Iteration can bound a small value:** The condition at \(n=2\) sharply restricts \(f(2)\).
\item **Integrality creates gaps:** The inequality \(f(3)^2>27\) becomes the stronger statement \(f(3)\ge6\).
\item **Compare the same number in different ways:** Multiplicativity translates numerical power comparisons into restrictions on function values.
\item **Use powers to amplify discrepancies:** A fixed multiplicative error cannot remain bounded through arbitrarily high powers.
\end{itemize}

\end{dlsummary}
